%%%%%%%%%%  scan idx 80  |  folio 69  %%%%%%%%%%

\underline{Corolario 2.18.} Un funtor
$F:\underline C_\Lambda\longrightarrow(\mathrm{Sets})$ es
prorepresentable si y solo si es homogéneo y
$\dim F(k[\varepsilon])<\infty$.

\underline{Corolario 2.19.} Todo cogrupoide liso en
$\widehat{\underline C}_\Lambda$ puede normalizarse salvo isomorfismo.

Dado un cogrupoide $(R_0,R_1,\ldots)$ en
$\widehat{\underline C}_\Lambda$, pongamos $\bar R_1=R_1/J$, donde $J$ es
el ideal de $R_1$ generado por
\[
 \{s_1(r)-s_2(r)\mid r\in R_0\}.
\]
Entonces $(R_0,\bar R_1,\ldots)$ define un esquema formal en grupos sobre
$R_0$ con la ley de composición inducida.

Sea ahora $\underline A$ un grupoide cofibrado sobre
$\underline C_\Lambda$, prorepresentado por un cogrupoide liso
normalizado $(R_0,R_1,\ldots)$ mediante $(\xi,\varphi)$. Sean
\[
 \underline C_k\lhook\joinrel\longrightarrow
 \underline C_{R_0}\lhook\joinrel\longrightarrow
 \underline C_{R_0\widehat\otimes_\Lambda R_0}
\]
las inclusiones correspondientes a las suprayecciones
\[
 R_0\widehat\otimes_\Lambda R_0\longrightarrow R_0\longrightarrow k,
\]
y sean $\operatorname{Aut}_\xi$ y $\operatorname{Aut}^{o}$ las
restricciones de $\operatorname{Isom}_\xi$ a
$\underline C_{R_0}$ y $\underline C_k$, respectivamente. Así,
$\operatorname{Aut}_\xi(S)$ es el grupo de automorfismos de $\xi(S)$ para
$S\in\operatorname{ob}(\underline C_{R_0})$, mientras que
$\operatorname{Aut}^{o}(S)$ es el grupo de automorfismos de $1_S$ para
$S\in\operatorname{ob}(\underline C_k)$, donde $1_S$ es la imagen
directa de $1\in\underline A(k)$ por $k\longrightarrow S$.

Como $R_1$ prorepresenta $\operatorname{Isom}_\xi$ sobre
$\underline C_{R_0\widehat\otimes_\Lambda R_0}$, se sigue que
$\bar R_1$ prorepresenta $\operatorname{Aut}_\xi$ sobre
$\underline C_{R_0}$, y que
$k\otimes_{R_0}\bar R_1$ prorepresenta $\operatorname{Aut}^{o}$ sobre
$\underline C_k$. Aunque $\underline A$ está prorepresentado por el
cogrupoide liso $(R_0,R_1,\ldots)$, ni
$(R_0,\bar R_1,\ldots)$ ni
$(k,k\otimes_{R_0}\bar R_1,\ldots)$ son lisos en general. Más
precisamente:

\underline{Corolario 2.20.} Sea $\underline A$ un
\sisi{cogrupoide homogéneo}{grupoide cofibrado homogéneo} sobre
$\underline C_\Lambda$, prorepresentado por un cogrupoide
liso normalizado $(R_0,R_1,\ldots)$ mediante $(\xi,\varphi)$.

\begin{enumerate}[label=(\alph*)]
\item $(R_0,\bar R_1,\ldots)$ y
$(k,k\otimes_{R_0}\bar R_1,\ldots)$ prorepresentan
$\operatorname{Aut}_\xi$ y $\operatorname{Aut}^{o}$, respectivamente.

\item Las afirmaciones siguientes son equivalentes:
\begin{enumerate}[label=(\roman*)]
\item el esquema formal en grupos
$\underline h_{(R_0,\bar R_1,\ldots)}$ es formalmente liso sobre $R_0$;
\item para toda flecha suprayectiva $a'\longrightarrow a$ en
$\underline A$, la aplicación
\[
 \operatorname{Aut}_{P(a')}(a')\longrightarrow
 \operatorname{Aut}_{P(a)}(a)
\]
es suprayectiva;
\item $R_0$ prorepresenta el funtor $[\underline A]$ sobre
$\underline C_\Lambda$.
\end{enumerate}
\end{enumerate}

%%%%%%%%%%  scan idx 81  |  folio 70  %%%%%%%%%%

\underline{Demostración.} Ya hemos señalado (a). Para (b), la
equivalencia (i)$\Leftrightarrow$(ii) es clara, mientras que
(ii)$\Leftrightarrow$(iii) se sigue de 2.7 y 2.15, puesto que
\[
 \widetilde\xi:\underline h_{R_0}\longrightarrow[\underline A]
\]
es mínimamente liso. \hfill$\square$

\underline{Observación 2.21.} Sea $\underline A$ un grupoide cofibrado
sobre $\underline C_\Lambda$, y sea
$\underline A|\underline C_k$ el cambio de base de
$p:\underline A\longrightarrow\underline C_\Lambda$ por medio de
$\underline C_k\lhook\joinrel\longrightarrow\underline C_\Lambda$.
Si $\underline A$ está prorepresentado por $(R_0,R_1,\ldots)$ mediante
$(\xi,\varphi)$, entonces $\underline A|\underline C_k$ está
prorepresentado por
\[
 (k\otimes_\Lambda R_0,k\otimes_\Lambda R_1,\ldots)
\]
mediante $(k\otimes_\Lambda\xi,k\otimes_\Lambda\varphi)$. Aquí
$k\otimes_\Lambda\xi$ es la imagen directa de $\xi$ por
$R_0\longrightarrow k\otimes_\Lambda R_0$, y
$k\otimes_\Lambda\varphi$ es la imagen de $\varphi$ por
\[
 \operatorname{Isom}_{\underline A(\sisi{R_0}{R_1})}
   \bigl(\xi(s_1),\xi(s_2)\bigr)
 \longrightarrow
 \operatorname{Isom}_{\underline A(k\otimes_\Lambda R_1)}
   \bigl(k\otimes_\Lambda\xi(s_1),
         k\otimes_\Lambda\xi(s_2)\bigr).
\]
Si $(R_0,R_1,\ldots)$ está normalizado, también lo está
$(k\otimes_\Lambda R_0,k\otimes_\Lambda R_1,\ldots)$; además, la
dimensión relativa de $R_1$ sobre $R_0$ es igual a la de
$k\otimes_\Lambda R_1$ sobre $k\otimes_\Lambda R_0$, puesto que $R_1$ es
liso sobre $R_0$.

\underline{Observación 2.22.} Supongamos que $\underline A$ está
prorepresentado por cogrupoides lisos $(R_0,R_1,\ldots)$ y
$(R'_0,R'_1,\ldots)$, con $(R_0,R_1,\ldots)$ normalizado. Obtenemos un
homomorfismo
\[
 (i_0,i_1):(R_0,R_1,\ldots)\longrightarrow(R'_0,R'_1,\ldots).
\]
Por 2.10(b), $R'_0$ es un anillo de series de potencias formales sobre
$R_0$, y $(i_0,i_1)$ induce un isomorfismo
\[
 R'_0\otimes_{R_0}(R_0,R_1,\ldots)
 \xrightarrow{\ \sim\ }(R'_0,R'_1,\ldots).
\]
